% from ega0-20-fr.tex
\oldpage[0\textsubscript{IV-1}]{129}
En outre, comme le diagramme
\[
\xymatrix{
  \mathrm{P}^1_{B/A} \ar[r]^{u'} & \mathrm{P}^1_{C/A} \\
  B \ar[u]^{p_2} \ar[r]^u & C \ar[u]_{p_2}
}
\tag{20.5.2.5}
\label{0.20.5.2.5-fr}
\]
est aussi commutatif, on en déduit que le diagramme
\[
\xymatrix{
  \Omega^1_{B/A} \ar[r]^{u''} & \Omega^1_{C/A} \\
  B \ar[u]^{d_{B/A}} \ar[r]^u & C \ar[u]_{d_{C/A}}
}
\tag{20.5.2.6}
\label{0.20.5.2.6-fr}
\]
est commutatif.

Enfin, si $w:C\to D$ est un second homomorphisme de $A$-algèbres, on a la
propriété de transitivité
\begin{equation}
\tag{20.5.2.7}
\label{0.20.5.2.7-fr}
  (w\circ u)_{D/B/A}
  =w_{D/C/A}\circ(u_{C/B/A}\otimes\mathrm{I}_D)
\end{equation}
comme il résulte de la définition.
\end{env}

\begin{env}[20.5.3]
\label{0.20.5.3-fr}
Soient maintenant $A$, $B$ deux anneaux, $v:A\to B$ un homomorphisme
d'anneaux, $C$ une $B$-algèbre qui devient une $A$-algèbre au moyen de $v$ ;
alors l'application canonique $v_0:C\otimes_A C\to C\otimes_B C$ est un
\emph{di-homomorphisme} surjectif d'algèbres (relatif à $v:A\to B$) tel que
le diagramme
\[
\xymatrix{
  C\otimes_A C \ar[r]^{v_0} \ar[d]_{p_{C/A}} &
  C\otimes_B C \ar[d]^{p_{C/B}} \\
  C \ar[r]^{\mathrm{I}_C} & C
}
\tag{20.5.3.1}
\label{0.20.5.3.1-fr}
\]
soit commutatif ; par passage aux quotients, on en déduit un
di-homomorphisme d'algèbres
\begin{equation}
\tag{20.5.3.2}
\label{0.20.5.3.2-fr}
  v':\mathrm{P}^1_{C/A}\to\mathrm{P}^1_{C/B}
\end{equation}
tel que le diagramme
\[
\xymatrix{
  \mathrm{P}^1_{C/A} \ar[r]^{v'} & \mathrm{P}^1_{C/B} \\
  C \ar[u]^{p_1} \ar[r]^{\mathrm{I}_C} & C \ar[u]_{p_1}
}
\]
soit commutatif. Comme $v_0$ applique $\mathfrak{J}_{C/A}$ dans
$\mathfrak{J}_{C/B}$ on obtient, en restreignant $v'$ à
$\Omega^1_{C/A}$, une application
\begin{equation}
\tag{20.5.3.3}
\label{0.20.5.3.3-fr}
  v_{C/B/A}:\Omega^1_{C/A}\to\Omega^1_{C/B}
\end{equation}
qui est un homomorphisme de $C$-modules.

